% Generated draft: requires mathematical and source review.
\documentclass[11pt]{article}
\usepackage[margin=25mm]{geometry}
\usepackage{fontspec}
\setmainfont{Latin Modern Roman}
\usepackage{amsmath,amssymb,mathtools,unicode-math}
\setmathfont{Latin Modern Math}
\usepackage{enumitem,needspace,xurl,hyperref}
\setlength{\parindent}{0pt}
\setlength{\parskip}{5pt}
\setlength{\emergencystretch}{3em}
\newcommand{\N}{\mathbb{N}}
\newcommand{\Z}{\mathbb{Z}}
\newcommand{\Q}{\mathbb{Q}}
\newcommand{\R}{\mathbb{R}}
\newcommand{\C}{\mathbb{C}}
\newcommand{\F}{\mathbb{F}}
\newcommand{\A}{\mathbb{A}}
\newcommand{\PP}{\mathbb{P}}
\newcommand{\E}{\mathbb{E}}
\newcommand{\eps}{\varepsilon}
\newcommand{\dd}{\,\mathrm{d}}
\begin{document}
\section*{UnsolvedMath ID: 30001388}
\textbf{Source problem code:} catalogue entry
\subsection{Definitions and mathematical statement}
The supplied record does not specify the class of maps \(f\). Consequently, its intended domain of quantification remains unresolved. The following makes the question precise under an explicitly conditional interpretation in which \(f\colon\C\to\C\) is transcendental entire; this hypothesis must be confirmed against the source before accepting the formulation.

Under this interpretation, let \(f^0=\operatorname{id}_{\C}\) and \(f^{n+1}=f\circ f^n\). Define the escaping set by
\[
I(f)=\{z\in\C: |f^n(z)|\longrightarrow\infty\text{ as }n\to\infty\}.
\]
The Fatou set \(F(f)\) consists of points having a neighborhood on which \(\{f^n:n\geq0\}\) is normal, with convergence measured in the spherical metric. A Fatou component is a connected component of \(F(f)\). Here a periodic escaping Fatou component, called a Baker domain in the supplied record, means a Fatou component \(U\) for which some integer \(p\geq1\) satisfies \(f^p(U)\subseteq U\) and every \(z\in U\) belongs to \(I(f)\).

The source question, under these hypotheses, is whether
\[
\forall f\colon\C\to\C\text{ transcendental entire},\quad
\forall U\text{ a periodic escaping Fatou component of }f,\quad
\partial_{\C}U\cap I(f)\neq\varnothing.
\]
Thus the boundary point sought is a finite complex point. If the intended class includes meromorphic maps or maps on another domain, the definitions and quantifiers require further specification.

The supplied literature assessment uses the wording that a Baker domain ``can have'' an escaping boundary point. Read literally, this is an existence question, whereas the displayed source statement asks whether every such component must have one. The formulation above follows the explicit original and clean statements; the assessment does not resolve this difference.
\subsection{Short English statement}
Must every periodic escaping Fatou component have an escaping point on its boundary? The displayed formulation uses an explicitly conditional transcendental entire setting because the supplied record does not specify the class of maps. The supplied open status is dated curation, not an independently verified conclusion.
\subsection{Sources}
\begin{itemize}
\item Supplied local UnsolvedMath registry record, \texttt{data/unsolvedmath/problems.json}, permanent ID \(30001388\), source problem code \texttt{catalogue entry}, title \emph{Escaping Boundary Points of Baker Domains}, in the Oberwolfach Reports Open Problems collection. No individual proposer is supplied; the proposed year is recorded as 2009.
\item Supplied primary citation: \emph{Mini-Workshop: The Escaping Set in Transcendental Dynamics} (2010), Oberwolfach Reports, pp.\ 2927--2964, DOI \texttt{10.4171/owr/2009/54}. \url{https://doi.org/10.4171/owr/2009/54}. This is also the literature-source link labeled ``original OWR report.'' The supplied record labels the statement \texttt{corrected\_verified} and reports comparison with the official report; that verification is supplied provenance, not a verification performed for this draft.
\item Supplied literature-source label: ``Bergweiler Rempe 2025 survey.'' \url{https://doi.org/10.1365/s13291-025-00300-1}. No further bibliographic details are supplied.
\item Literature checked at: \texttt{2026-08-21}. This is supplied dated curation, not independently verified. The supplied status is \emph{open}; its assessment says that the general boundary question remains open and that later boundary results require additional hypotheses or special classes. The supplied background explicitly requests independent verification of time-sensitive status and cited claims. Its ``can have'' wording must be distinguished from the universal question in the source statement.
\end{itemize}
\end{document}
